% GENERATED FILE. Edit canonical lessons, then run npm run generate:books.
% canonical-source-hash: fnv1a64:c93ac462ba564add
% canonical-lessons: PT-C511-torneio, PT-C511-taca, PT-C511-medalha, PT-C511-podio, PT-C511-recorde

\chapter{Taças e medalhas}
\label{ch:pt-tacas-e-medalhas}
\hlchaptermodality{511}

\begin{chapteropening}
\textbf{By the end of this chapter:} \emph{I can talk about tournaments, cups and medals.}

Complete the last lesson of chapter 511: I can talk about tournaments, cups and medals.
\end{chapteropening}

\section[o torneio]{o torneio — tournament}

\label{lesson:PT-C511-torneio}



\begin{quote}
Before the new one: say the Portuguese for to defeat, then the Portuguese for to knock out.
\end{quote}

\subsection*{You'll want to know: o torneio}
\textbf{o torneio} — \textquotedblleft{}tournament\textquotedblright{}.

\begin{quote}
A escola organiza um torneio de ténis.
\end{quote}

\begin{quote}
\emph{The school is organising a tennis tournament.}
\end{quote}

Medieval knights fought in \emph{torneios} too.

\begin{cousinweb}
\emph{torneio} came into Old Portuguese as \emph{torneyo} from Old Occitan \emph{tornei} (EtymDB), which goes back to Latin \emph{tornāre} \textquotedblleft{}to turn\textquotedblright{}.
\end{cousinweb}

\begin{grammarlens}[title={What you've built}]
One more word for telling, explaining and describing at B1.
\end{grammarlens}

\subsection*{Your turn}
\begin{itemize}
\raggedright
  \item \emph{Say it:} \emph{o torneio}
  \item \emph{Say it:} \emph{o torneio}, once more
  \item \emph{Say it:} a sentence of your own with \emph{o torneio}
  \item \emph{Recall:} say the Portuguese for to defeat, then the Portuguese for to knock out, then say \emph{o torneio} again
\end{itemize}

\begin{tcolorbox}[breakable,colback=teal!4,colframe=teal!35!black,title={Before you move on}]
What does o torneio mean? (\textquotedblleft{}Tournament\textquotedblright{}.) Say it once more.
\end{tcolorbox}
\section[a taça]{a taça — cup, trophy; bowl}

\label{lesson:PT-C511-taca}



\begin{quote}
Before the new one: say the Portuguese for to knock out, then the Portuguese for tournament.
\end{quote}

\subsection*{You'll want to know: a taça}
\textbf{a taça} — \textquotedblleft{}cup, trophy; bowl\textquotedblright{}.

\begin{quote}
A equipa levantou a taça no fim do jogo.
\end{quote}

\begin{quote}
\emph{The team lifted the cup at the end of the match.}
\end{quote}

Also a wine glass or a bowl: \emph{uma taça de fruta}.

\begin{cousinweb}
\emph{taça} came through Spanish \emph{taza} from Arabic \emph{ṭāsa} \textquotedblleft{}a bowl\textquotedblright{} (EtymDB), which goes back to Persian.
\end{cousinweb}

\begin{grammarlens}[title={What you've built}]
One more word for telling, explaining and describing at B1.
\end{grammarlens}

\subsection*{Your turn}
\begin{itemize}
\raggedright
  \item \emph{Say it:} \emph{a taça}
  \item \emph{Say it:} \emph{a taça}, once more
  \item \emph{Say it:} a sentence of your own with \emph{a taça}
  \item \emph{Recall:} say the Portuguese for to knock out, then the Portuguese for tournament, then say \emph{a taça} again
\end{itemize}

\begin{tcolorbox}[breakable,colback=teal!4,colframe=teal!35!black,title={Before you move on}]
What does a taça mean? (\textquotedblleft{}Cup, trophy; bowl\textquotedblright{}.) Say it once more.
\end{tcolorbox}
\section[a medalha]{a medalha — medal}

\label{lesson:PT-C511-medalha}



\begin{quote}
Before the new one: say the Portuguese for tournament, then the Portuguese for cup.
\end{quote}

\subsection*{You'll want to know: a medalha}
\textbf{a medalha} — \textquotedblleft{}medal\textquotedblright{}.

\begin{quote}
Ganhou uma medalha de ouro.
\end{quote}

\begin{quote}
\emph{She won a gold medal.}
\end{quote}

\emph{Medalha de prata}, \emph{de bronze} — silver, bronze.

\begin{cousinweb}
\emph{medalha} comes from Italian \emph{medaglia} (EtymDB), which goes back to Latin \emph{metallum} \textquotedblleft{}metal\textquotedblright{}, from Greek.
\end{cousinweb}

\begin{grammarlens}[title={What you've built}]
One more word for telling, explaining and describing at B1.
\end{grammarlens}

\subsection*{Your turn}
\begin{itemize}
\raggedright
  \item \emph{Say it:} \emph{a medalha}
  \item \emph{Say it:} \emph{a medalha}, once more
  \item \emph{Say it:} a sentence of your own with \emph{a medalha}
  \item \emph{Recall:} say the Portuguese for tournament, then the Portuguese for cup, then say \emph{a medalha} again
\end{itemize}

\begin{tcolorbox}[breakable,colback=teal!4,colframe=teal!35!black,title={Before you move on}]
What does a medalha mean? (\textquotedblleft{}Medal\textquotedblright{}.) Say it once more.
\end{tcolorbox}
\section[o pódio]{o pódio — podium}

\label{lesson:PT-C511-podio}



\begin{quote}
Before the new one: say the Portuguese for cup, then the Portuguese for medal.
\end{quote}

\subsection*{You'll want to know: o pódio}
\textbf{o pódio} — \textquotedblleft{}podium\textquotedblright{}.

\begin{quote}
Três atletas subiram ao pódio.
\end{quote}

\begin{quote}
\emph{Three athletes went up on the podium.}
\end{quote}

\emph{Subir ao pódio} — to finish in the top three.

\begin{cousinweb}
\emph{pódio} is a learned borrowing from Latin \emph{podium} (EtymDB), from Greek \emph{pódion} \textquotedblleft{}a base\textquotedblright{}, built on \emph{poús} \textquotedblleft{}foot\textquotedblright{}.
\end{cousinweb}

\begin{grammarlens}[title={What you've built}]
One more word for telling, explaining and describing at B1.
\end{grammarlens}

\subsection*{Your turn}
\begin{itemize}
\raggedright
  \item \emph{Say it:} \emph{o pódio}
  \item \emph{Say it:} \emph{o pódio}, once more
  \item \emph{Say it:} a sentence of your own with \emph{o pódio}
  \item \emph{Recall:} say the Portuguese for cup, then the Portuguese for medal, then say \emph{o pódio} again
\end{itemize}

\begin{tcolorbox}[breakable,colback=teal!4,colframe=teal!35!black,title={Before you move on}]
What does o pódio mean? (\textquotedblleft{}Podium\textquotedblright{}.) Say it once more.
\end{tcolorbox}
\section[o recorde]{o recorde — record (in sport)}

\label{lesson:PT-C511-recorde}



\begin{quote}
Before the new one: say the Portuguese for medal, then the Portuguese for podium.
\end{quote}

\subsection*{You'll want to know: o recorde}
\textbf{o recorde} — \textquotedblleft{}record (in sport)\textquotedblright{}.

\begin{quote}
Bateu o recorde nacional dos cem metros.
\end{quote}

\begin{quote}
\emph{He broke the national hundred-metre record.}
\end{quote}

\emph{Bater um recorde} — to break a record.

\begin{cousinweb}
\emph{recorde} was borrowed from English \emph{record} (EtymDB), which came from Old French.
\end{cousinweb}

\begin{grammarlens}[title={What you've built}]
One more word for telling, explaining and describing at B1.
\end{grammarlens}

\subsection*{Your turn}
\begin{itemize}
\raggedright
  \item \emph{Say it:} \emph{o recorde}
  \item \emph{Say it:} \emph{o recorde}, once more
  \item \emph{Say it:} a sentence of your own with \emph{o recorde}
  \item \emph{Recall:} say the Portuguese for medal, then the Portuguese for podium, then say \emph{o recorde} again
\end{itemize}

\begin{tcolorbox}[breakable,colback=teal!4,colframe=teal!35!black,title={Before you move on}]
What does o recorde mean? (\textquotedblleft{}Record (in sport)\textquotedblright{}.) Say it once more.
\end{tcolorbox}
